Mapa průměrných mezd v~sektoru informačních a~komunikačních technologií (\ac{NACE} J) v~EU27 ukazuje, že ČR dosahuje relativně lepší pozice než v~ostatních sledovaných odvětvích --- Praha je uznávaným ICT centrem s~mezinárodně konkurenceschopnými mzdami v~tomto sektoru.
Nicméně i~v~ICT je pozice ČR za AT nebo DE i~po kontrole cenové hladiny: přepočteno na \ac{PPS} je průměrná ICT mzda v~ČR přibližně o~\SIrange{15}{20}{\percent} nižší než v~Německu, přičemž IT odborníci tuto mezeru jednoznačně vnímají a~reagují emigrací nebo remote prací pro zahraniční klienty.
Paradoxem ICT je, že jde o~sektor s~nejsilnější individuální vyjednávací pozicí (vysoká poptávka, nízká zastupitelnost, globální mobilita) --- přesto i~zde absence odvětvové \ac{KS} umožňuje zaměstnavatelům udržet mzdy pod potenciálem, zejména pro méně zkušené pracovníky bez přímé mezinárodní alternativy.
Doporučení: odvětvová \ac{KS} v~ICT sektoru by nemusela omezovat tržní prémie pro seniorní odborníky (ti si tyto prémie vyjednají individuálně), ale zajistila by minimální tarifní ochranu pro juniorní a~mid-level pracovníky, kteří jsou v~informační asymetrii vůči zaměstnavateli.

Mapa průměrných mezd v~sektoru informačních a~komunikačních technologií (\ac{NACE} J) v~EU27 ukazuje, že ČR dosahuje relativně lepší pozice než v~ostatních sledovaných odvětvích --- Praha je uznávaným ICT centrem s~mezinárodně konkurenceschopnými mzdami v~tomto sektoru.
Nicméně i~v~ICT je pozice ČR za AT nebo DE i~po kontrole cenové hladiny: přepočteno na \ac{PPS} je průměrná ICT mzda v~ČR přibližně o~\SIrange{15}{20}{\percent} nižší než v~Německu, přičemž IT odborníci tuto mezeru jednoznačně vnímají a~reagují emigrací nebo remote prací pro zahraniční klienty.
Paradoxem ICT je, že jde o~sektor s~nejsilnější individuální vyjednávací pozicí (vysoká poptávka, nízká zastupitelnost, globální mobilita) --- přesto i~zde absence odvětvové \ac{KS} umožňuje zaměstnavatelům udržet mzdy pod potenciálem, zejména pro méně zkušené pracovníky bez přímé mezinárodní alternativy.
Doporučení: odvětvová \ac{KS} v~ICT sektoru by nemusela omezovat tržní prémie pro seniorní odborníky (ti si tyto prémie vyjednají individuálně), ale zajistila by minimální tarifní ochranu pro juniorní a~mid-level pracovníky, kteří jsou v~informační asymetrii vůči zaměstnavateli.

Mapa průměrných mezd v~sektoru informačních a~komunikačních technologií (\ac{NACE} J) v~EU27 ukazuje, že ČR dosahuje relativně lepší pozice než v~ostatních sledovaných odvětvích --- Praha je uznávaným ICT centrem s~mezinárodně konkurenceschopnými mzdami v~tomto sektoru.
Nicméně i~v~ICT je pozice ČR za AT nebo DE i~po kontrole cenové hladiny: přepočteno na \ac{PPS} je průměrná ICT mzda v~ČR přibližně o~\SIrange{15}{20}{\percent} nižší než v~Německu, přičemž IT odborníci tuto mezeru jednoznačně vnímají a~reagují emigrací nebo remote prací pro zahraniční klienty.
Paradoxem ICT je, že jde o~sektor s~nejsilnější individuální vyjednávací pozicí (vysoká poptávka, nízká zastupitelnost, globální mobilita) --- přesto i~zde absence odvětvové \ac{KS} umožňuje zaměstnavatelům udržet mzdy pod potenciálem, zejména pro méně zkušené pracovníky bez přímé mezinárodní alternativy.
Doporučení: odvětvová \ac{KS} v~ICT sektoru by nemusela omezovat tržní prémie pro seniorní odborníky (ti si tyto prémie vyjednají individuálně), ale zajistila by minimální tarifní ochranu pro juniorní a~mid-level pracovníky, kteří jsou v~informační asymetrii vůči zaměstnavateli.

Mapa průměrných mezd v~sektoru informačních a~komunikačních technologií (\ac{NACE} J) v~EU27 ukazuje, že ČR dosahuje relativně lepší pozice než v~ostatních sledovaných odvětvích --- Praha je uznávaným ICT centrem s~mezinárodně konkurenceschopnými mzdami v~tomto sektoru.
Nicméně i~v~ICT je pozice ČR za AT nebo DE i~po kontrole cenové hladiny: přepočteno na \ac{PPS} je průměrná ICT mzda v~ČR přibližně o~\SIrange{15}{20}{\percent} nižší než v~Německu, přičemž IT odborníci tuto mezeru jednoznačně vnímají a~reagují emigrací nebo remote prací pro zahraniční klienty.
Paradoxem ICT je, že jde o~sektor s~nejsilnější individuální vyjednávací pozicí (vysoká poptávka, nízká zastupitelnost, globální mobilita) --- přesto i~zde absence odvětvové \ac{KS} umožňuje zaměstnavatelům udržet mzdy pod potenciálem, zejména pro méně zkušené pracovníky bez přímé mezinárodní alternativy.
Doporučení: odvětvová \ac{KS} v~ICT sektoru by nemusela omezovat tržní prémie pro seniorní odborníky (ti si tyto prémie vyjednají individuálně), ale zajistila by minimální tarifní ochranu pro juniorní a~mid-level pracovníky, kteří jsou v~informační asymetrii vůči zaměstnavateli.

\input{texparts/python/eu_odvetvove_mzdy_mapa_J}



